\chapter{Evaluation and Reflection}
\label{ch:evaluation}

\section{Designing an evaluation when there is no ground truth}
\label{sec:eval-design}

The obvious evaluation is accuracy against real outcomes, and it is mostly
unavailable. Fusion21's own Doc 4 states that the overwhelming majority of
procurement disputes settle before proceedings, so the cases that produce a
published judgment are a biased minority; of the 23 in this corpus, nine were
decided on an interim application where the underlying merits were never reached
at all. Any single accuracy number computed over that set would be both small and
misleading.

The evaluation is therefore built as a stack of partial checks
(Figure~\ref{fig:evidence}), each answering a question the others cannot, and each
reported with the sample size that produced it. Every figure below comes from an
analysis script in the repository run over committed logs, and every script was
re-run immediately before submission. The corpus at that point was 299 completed
negotiations across 59 batches.

\begin{figure}[htbp]
\centering
\begin{tikzpicture}[font=\footnotesize,
  band/.style={draw=f21ink,rounded corners=2pt,line width=0.7pt,align=left,
               inner sep=5pt,text width=11.6cm}]
  \node[band,fill=f21mist] (a)
    {\textbf{Does it produce well-formed output at all?}\quad
     Structural compliance over 2{,}865 structured responses~(\S\ref{sec:structural})};
  \node[band,below=2.6mm of a,fill=f21mist] (b)
    {\textbf{Does it agree with reality where reality is knowable?}\quad
     14 merits dispositions~(\S\ref{sec:agreement})};
  \node[band,below=2.6mm of b,fill=f21green!10] (c)
    {\textbf{Does the architecture earn its complexity?}\quad
     Three baselines, one ablation, $n=6$ and $n=18$~(\S\ref{sec:baselines})};
  \node[band,below=2.6mm of c,fill=f21green!10] (d)
    {\textbf{Does a prompt change alter judgement or only confidence?}\quad
     Controlled V3/V4 ablation, $n=8$ per cell, Fisher's exact~(\S\ref{sec:ablation})};
  \node[band,below=2.6mm of d,fill=f21amber!10] (e)
    {\textbf{Is the reasoning grounded in what the model was given?}\quad
     Three independent fabrication screens~(\S\ref{sec:fabrication})};
  \node[band,below=2.6mm of e,fill=f21amber!10] (f)
    {\textbf{Is the output usable by the person who has to act on it?}\quad
     Concession dynamics, BATNA self-report, readability~(\S\ref{sec:dynamics}--\S\ref{sec:readability})};
\end{tikzpicture}
\caption[Evaluation design]{The evaluation as a stack of partial checks. No single
row is sufficient; the fabrication screens in particular exist because a system can
score well on every row above them while inventing the evidence it reasons from.}
\label{fig:evidence}
\end{figure}

\section{The corpus}
\label{sec:corpus}

Twenty-three real UK procurement judgments were researched, individually verified
against BAILII or the National Archives case law service, and written into the
project as dense source texts before being run through the same extraction path a
user upload takes. Seven further candidates were investigated and rejected during
that process, two of which turned out on closer reading not to be procurement
cases at all. Table~\ref{tab:corpus-summary} summarises the composition; the full
list with citations is in Appendix~\ref{app:corpus}.

\begin{table}[htbp]
\centering
\footnotesize
\caption[Corpus composition]{Composition of the 23-case corpus. Only the merits
group supports an agreement claim; the interim group is reported for completeness
and excluded from every accuracy figure in this chapter.}
\label{tab:corpus-summary}
\begin{tabular}{@{}lccp{5.6cm}@{}}
\toprule
\textbf{Group} & \textbf{Cases} & \textbf{CA won / lost} & \textbf{Why grouped this way} \\
\midrule
Merits disposition & 14 & 6 / 8 &
  A court ruled on the substantive question this system assesses \\
Interim only & 9 & --- &
  Suspension or summary-judgment rulings; the merits were never decided, so
  comparing a merits prediction to them is not a check \\
\bottomrule
\end{tabular}
\end{table}

The corpus grew in three stages, and the constraint on growth changed as it did.
While inference ran on a CPU-bound local model, a single negotiation took around
seventeen minutes, which made a twenty-case corpus a compute problem. Once
inference moved to a tunnelled A10G GPU, per-call latency fell roughly thirty-fold
and the bottleneck moved to research and verification: the largest expansion added
fifteen cases in one session, with the compute for all fifteen finishing inside an
hour and the searching and checking taking the rest of it.

\section{Structural reliability}
\label{sec:structural}

Across the 50 batches carrying the instrumentation described in
\S\ref{sec:compliance-method}, 2{,}861 of 2{,}865
structured responses required no repair, giving a corpus-wide structural
compliance rate of \textbf{0.9986}. Forty-eight of the 50 batches were clean at
1.00; the worst single batch was 0.906 and predates the repair-retry pattern
described below.

That pattern is itself a result. When the anti-fabrication and
governing-legislation instructions were added to the pre-negotiation prompts, the
rate of the model dropping a required JSON field rose measurably, first observed as
three failures in an eight-run batch where the entire prior corpus had none. The
cause was not the content of the instruction but its length: a longer,
instruction-denser prompt degraded adherence to the output-format instruction
sitting at the end of it. The fix was a single repair retry that re-sends the same
message with the validation error appended, the same pattern the extraction agent
already used. This interaction recurred, and is worth naming for anyone hardening
prompts against a small model: \emph{adding an instruction has a cost paid
somewhere else in the same prompt}, and the cost is not visible unless something
counts it.

\section{Agreement with real dispositions}
\label{sec:agreement}

On the 14 cases with a genuine merits disposition, the simulated outcome direction
matches the real one in \textbf{13}. The exception is Turning Point v Norfolk
County Council, and it is not a bug. Reading the transcript rather than the
outcome distribution, the simulated Court's reasoning is almost verbatim the real
claimant's losing argument: that the authority should have sought clarification of
a caveated tender before rejecting it. The real judge rejected exactly that
argument, holding the no-caveats rule fair and commonly used. Nothing was invented
and the reasoning engages the real facts throughout; the model simply reached a
different legal conclusion on a genuinely arguable point. That is a category of
error worth separating from fabrication, and Section~\ref{sec:fabrication} does so.

Two caveats attach to the 13. On Faraday v West Berkshire the direction is correct
but the reasoning is not trustworthy, for reasons given below. On Consultant
Connect the direction is correct but the remedy shape is wrong: the real court
imposed financial penalties on three NHS bodies, and this system's outcome
vocabulary has no way to express that, so ``re-evaluation'' is the closest
available match to ``the authority was at fault''. The same gap appears on Faraday,
whose real remedy was a declaration of ineffectiveness, and on Woods, where the
court did not merely order the scoring redone but determined the corrected winner
itself. The vocabulary was built for scoring challenges and shows it.

Across all 299 runs the outcome distribution is heavily skewed: re-evaluation
80.6\%, no remedy 10.4\%, deadlock 5.7\%, disclosure ordered 3.0\%. One run
produced ``request for additional information'', a string outside the closed
vocabulary entirely --- a single instance in 299, from an early four-run batch, and
precisely the drift the downstream module raises on rather than silently averaging
away. Median resolution takes one round; the maximum observed is five.

\section{Baselines and ablations: does the architecture earn itself?}
\label{sec:baselines}

This is the question an examiner asks first, so it was tested directly, on the six
cases where a merits disposition existed at the time of the comparison. Three
baselines were built: a majority-class heuristic with no model call at all, a
single zero-shot prompt deliberately \emph{not} hardened the way the Court prompt
is, and an ablation of the full pipeline with the Court node replaced by one that
cannot resolve.

\begin{table}[htbp]
\centering
\footnotesize
\caption[Baseline comparison]{Directional accuracy against real dispositions, and
the no-Court ablation. The imbalance matters: five of the six cases are ``CA
lost'', so always guessing the majority class scores 5/6 with no reasoning
whatsoever.}
\label{tab:baselines}
\begin{tabular}{@{}lccl@{}}
\toprule
\textbf{System} & \textbf{Runs} & \textbf{Direction correct} & \textbf{What it establishes} \\
\midrule
Majority-class heuristic & --- & 5 / 6 & The floor. No model call, no reasoning \\
Zero-shot single prompt & 5 per case & 6 / 6 & A naive prompt is not worse here \\
Full pipeline (V4) & 8 per case & 6 / 6 & Ties the naive prompt \\
No-Court ablation & 3 per case & --- & \textbf{0 of 18 runs resolved} \\
\bottomrule
\end{tabular}
\end{table}

The headline is a negative result and is reported as measured.
\textbf{At $n=6$ the full multi-agent pipeline and the naive single-prompt
baseline tie: both get every case right, and the gap to a predictor with no
reasoning at all is one case wide.} Any claim in this project about the
architecture's value cannot rest on accuracy. A second check tested whether the
unhardened prompt at least fabricated more, applying the same grounding screen
used on the Court agent to all 30 zero-shot reasoning texts: zero ungrounded
numbers were found. That is a genuine null result on this evidence, not a finding
massaged into a narrative.

What the ablation shows is different in kind. Removing the Court agent produces
deadlock in \textbf{all 18 runs, uniformly on every case}. This is not a trend
subject to sampling luck; it is architecturally necessary, because nothing else in
the graph can decide the parties have converged. A version of the ablation that
inferred resolution from conciliatory language in the dialogue would have defeated
the point, which is precisely that nothing here resolves without adjudication.
The multi-agent design's defensible claim is therefore about mechanism rather than
accuracy: the pipeline produces a modelled negotiation with declared interests,
BATNAs, rounds of concession, and a court reviewing the actual exchange, none of
which exists in a single call, and all of which the downstream components depend
on structurally.

\section{The V3/V4 prompt ablation}
\label{sec:ablation}

Replacing the Court agent's three-line guiding-principles block with Doc 1's six
judiciary interest categories is a single-variable change, and it has a consistent
effect. Across every case where V3 deadlocks at all, V4 reaches zero deadlock and
resolves in roughly half the rounds (Figure~\ref{fig:ablation},
Table~\ref{tab:ablation}).

\begin{figure}[htbp]
\centering
\begin{tikzpicture}[font=\footnotesize]
  % axes
  \draw[f21grey,line width=0.5pt] (0,0) -- (11.6,0);
  \foreach \y/\lab in {0/0, 1.4/.25, 2.8/.50, 4.2/.75, 5.6/1.0} {
    \draw[f21grey!45,line width=0.4pt] (0,\y) -- (11.6,\y);
    \node[anchor=east,f21grey,font=\scriptsize] at (-0.1,\y) {\lab};
  }
  \node[rotate=90,anchor=south,f21grey,font=\scriptsize] at (-0.95,2.8) {resolution rate};

  % five case groups; x centres
  \foreach \i/\lbl/\vthree in {0/Lancashire/4.9, 1/Faraday/4.2, 2/Parkingeye/4.9,
                               3/Alstom/2.8, 4/Woods/5.6} {
    \pgfmathsetmacro\xc{0.75 + \i*2.25}
    \fill[f21grey!55] ({\xc-0.62},0) rectangle ({\xc-0.06},\vthree);
    \fill[f21green]   ({\xc+0.06},0) rectangle ({\xc+0.62},5.6);
    \node[anchor=north,font=\scriptsize] at (\xc,-0.12) {\lbl};
  }
  % legend
  \fill[f21grey!55] (8.9,6.25) rectangle (9.25,6.55);
  \node[anchor=west,font=\scriptsize] at (9.32,6.4) {V3};
  \fill[f21green] (10.2,6.25) rectangle (10.55,6.55);
  \node[anchor=west,font=\scriptsize] at (10.62,6.4) {V4};
\end{tikzpicture}
\vspace{4mm}
\caption[V3 vs V4 resolution rates]{Resolution rate by Court prompt version,
$n=8$ per bar, five real cases. V4 sits at the ceiling everywhere; V3 varies by
case. Reaching 1.00 on every case is not automatically the better result --- see
the text.}
\label{fig:ablation}
\end{figure}

\begin{table}[htbp]
\centering
\footnotesize
\caption[Ablation significance]{V3/V4 resolution rates with Fisher's exact tests
\cite{fisher1922interpretation}.
No case reaches $p<0.05$. At $n=8$ per arm this test has low power, so these are
``not distinguishable from noise at this sample size'', not ``no difference''.}
\label{tab:ablation}
\begin{tabular}{@{}lcccc@{}}
\toprule
\textbf{Case} & \textbf{V3} & \textbf{V4} & \textbf{V3 avg.\ rounds} & \textbf{Fisher $p$} \\
\midrule
Lancashire Care & 0.875 & 1.000 & 2.75 & 1.000 \\
Faraday & 0.750 & 1.000 & 3.00 & 0.467 \\
Parkingeye & 0.875 & 1.000 & 1.50 & 1.000 \\
Alstom & 0.500 & 1.000 & 4.38 & 0.077 \\
Woods & 1.000 & 1.000 & 1.00 & 1.000 \\
\bottomrule
\end{tabular}
\end{table}

Two things follow, and they pull in opposite directions. The pattern is real and
consistent across four independent cases, and it matches the mechanism you would
predict from the change: Doc 1's fifth judicial category is efficient use of
judicial resources, which favours settlement and narrowing of issues. But nothing
here is statistically significant, and a prompt that resolves everything is not
obviously the more correct one. On Alstom, V4's manifest-error detection rate is
identical to V3's at 0.50 --- so the difference is entirely in what happens
\emph{after} a finding, not in whether one is made. The honest reading is that V4
changes remedy efficiency rather than judgement, which is what its explicit
precedence note was written to ensure, but ``V4 improves resolution rate'' should
not be cited as an unqualified positive.

\section{Fabrication: three distinct failure modes}
\label{sec:fabrication}

This is the project's central concern and the screens are correspondingly
paranoid. Three independent checks run over the corpus, summarised in
Table~\ref{tab:fabrication}; the failure modes they separate are shown in
Figure~\ref{fig:failures}.

\begin{table}[htbp]
\centering
\footnotesize
\caption[Fabrication screens]{The three grounding screens, over 299 logs.}
\label{tab:fabrication}
\begin{tabular}{@{}lcp{6.2cm}@{}}
\toprule
\textbf{Screen} & \textbf{Result} & \textbf{Reading} \\
\midrule
Numeric grounding (Step 2B) & 9 flagged of 425 (2.1\%) &
  All nine re-read individually; \textbf{none is a Court-originated fabrication} \\
Legal citation, regime & 162 of 702 (23.1\%) wrong regime &
  PCR 2015 or the EU Directive cited in a Procurement Act 2023 world \\
Legal citation, s.12 content & 7 of 21 (33.3\%) misattributed &
  Subsection contents swapped, e.g.\ transparency assigned to s.12(1)(a) \\
\bottomrule
\end{tabular}
\end{table}

The numeric screen extracts every score-shaped number from a qualitative
compliance assessment and checks it appears in the scenario description. Nine of
425 qualitative assessments flagged, and reading all nine changes the
interpretation entirely. None is the Court inventing a figure. Six reproduce a
pattern first seen on Alstom, where the negotiating agents invent quantitative
specificity and the Court declines to treat it as verified. Two are the Court
quoting a number attributively --- ``the Contracting Authority's position states
that the bid scored 92 out of 100'' --- and one is a Court repeating a party's own
negotiating proposal. The clearest of the nine is worth quoting, because the
detector fired on an instance of the discipline working exactly as designed:

\begin{quote}\itshape\small
``I do not have enough information to perform a calculation myself, as no specific
numbers or formula are provided in the scenario.''
\hfill --- Court agent, Alstom, round 2
\end{quote}

The screen's real signal is therefore upstream. The negotiating agents fabricate
factual premises, and the Court adopts them. The clearest instance is Faraday,
which is not a scoring dispute at all: the authority structured a
\pounds125\,m regeneration deal so as to avoid running a competition, and Faraday
never submitted a bid to be scored. Extraction captured this correctly, labelling
the dispute \texttt{process\_avoidance}. The negotiating agents did not preserve
it. In 3 of 8 V3 runs and 6 of 8 V4 runs, the bidder asserts a ``scoring
discrepancy between our bid and that of St Modwen Developments Ltd'' --- a
comparison that cannot exist, since neither party was ever scored. The cause is
traceable: every CA and bidder prompt in this system was developed against
scored-bid disputes, and neither contains a branch for a case where no bid was
ever scored. Confronted with an unfamiliar dispute type, both agents default to
the only script they have.

\begin{figure}[htbp]
\centering
\begin{tikzpicture}[font=\footnotesize,
  m/.style={draw,rounded corners=2pt,line width=0.9pt,align=left,inner sep=6pt,
            text width=3.55cm,minimum height=3.0cm,anchor=north west}]
  \node[m,draw=f21green,fill=f21green!8] at (0,0)
    {\textbf{1. Numeric fabrication}\\[2pt]
     \scriptsize Inventing a sub-score or formula the case does not state.\\[3pt]
     \textbf{Status:} addressed. Step 1 gating; 0 of 9 flagged instances
     Court-originated.};
  \node[m,draw=f21amber,fill=f21amber!10] at (4.05,0)
    {\textbf{2. Premise fabrication}\\[2pt]
     \scriptsize A negotiating agent inventing a fact about the dispute, which the
     Court then adopts.\\[3pt]
     \textbf{Status:} open. 6 of 8 runs on Faraday.};
  \node[m,draw=f21grey,fill=f21mist] at (8.1,0)
    {\textbf{3. Judgment divergence}\\[2pt]
     \scriptsize Grounded reasoning reaching a different legal conclusion than the
     real judge.\\[3pt]
     \textbf{Status:} not a defect. Turning Point.};
\end{tikzpicture}
\caption[Failure mode taxonomy]{Three failure modes that a single ``hallucination''
label would conflate. Hardening the adjudicator addresses the first and does
nothing about the second, because the second enters the system before the
adjudicator sees it.}
\label{fig:failures}
\end{figure}

\subsection{Legal citation}
\label{sec:citation}

Citation is the same discipline in a different field, and it holds far less well.
Of 702 citations logged, 162 (23.1\%) name the wrong regime entirely, citing the
Public Contracts Regulations 2015 or the EU Directive in a system explicitly
grounded in the Procurement Act 2023; one names a ``Public Contracts Regulations
2020'' that does not exist. Of the 21 citations specific enough to check against
section 12, 7 misattribute the subsection's contents. Both checks were verified
against legislation.gov.uk rather than assumed. The remaining 519 citations are
bucketed as not specific enough to check rather than silently passed as correct,
which is a deliberately conservative choice: it means the true error rate is
bounded below by these figures, not estimated by them.

This is the clearest limitation of prompt-only grounding. The citation-discipline
instruction reduced fabricated pinpoint citations, but it cannot supply the model
with knowledge of which regime governs a 2018 procurement, and an 8B model's
priors are saturated with pre-2023 material. Retrieval over the actual legislation
is the right fix and is the strongest argument for building the layer that was
scoped out.

\section{Negotiation dynamics and BATNA}
\label{sec:dynamics}

Two behavioural checks confirm the negotiation is a negotiation rather than two
monologues. Concession rate by the contracting authority rises from 0.013 in round
one to 0.470 in later rounds: it essentially never concedes on opening, and
concedes about half the time once the exchange is underway. The bidder concedes
far less, 0.000 to 0.075, which is consistent with what a bidder concession
actually means --- dropping the challenge, a much larger move than offering extra
transparency. Consecutive same-role message similarity averages 0.097 over 426
pairs, so the repetition suppression is working. Zero explicit retractions were
found, using a check narrowed after an earlier version misread ``willing to
withdraw our claim'' --- a settlement offer --- as a retraction.

Table~\ref{tab:batna} shows how each agent judges its own outcome against its
declared BATNA. The asymmetry is structural and matches the law: an authority that
successfully defends its award loses nothing, while a bidder that obtains a
re-evaluation has still spent the bid cost and may lose again.

\begin{table}[htbp]
\centering
\footnotesize
\caption[BATNA self-assessment]{Self-reported outcome relative to declared BATNA,
$n=299$ per role, classified by a documented keyword heuristic. The heuristic is
not a ground-truth label, and the large ``unclear'' fractions are reported rather
than redistributed.}
\label{tab:batna}
\begin{tabular}{@{}lcccc@{}}
\toprule
\textbf{Role} & \textbf{Beats} & \textbf{Matches} & \textbf{Falls short} & \textbf{Unclear} \\
\midrule
Contracting Authority & 71\% & 1\% & 2\% & 27\% \\
Aggrieved Bidder & 24\% & 0\% & 21\% & 56\% \\
\bottomrule
\end{tabular}
\end{table}

An earlier run of this analysis over 63 logs reported the authority as never
falling short. At 299 logs it falls short in five cases, all of which read as
genuine self-reported shortfalls on inspection. The absolute claim was never safe
to make from the smaller sample, and the correction is a caution about how any
proportion in this project should be cited before $N$ is much larger.

\section{Readability: the summariser does not meet its own claim}
\label{sec:readability}

The Summary agent's schema describes its output as ``a 3--4 sentence summary a
non-lawyer could understand''. Measured over all 299 summaries, mean Flesch
Reading Ease \cite{flesch1948readability} is \textbf{24.7} (median 25.9), which is
the ``very difficult, graduate level'' band; mean Flesch--Kincaid grade level
\cite{kincaid1975derivation} is 14.7. The best single
summary in the corpus scores 57.5, just short of the 60--70 plain-English band, and
the worst is negative. The figure is essentially unchanged from an earlier run over
a different and smaller corpus, so it is a stable property of the prompt rather
than a sampling artefact. The agent that exists specifically to make the output
accessible does not, on a standard measure, succeed --- and it would not have been
noticed without measuring it, because the summaries read fluently.

\section{The risk screen against real cases}
\label{sec:risk-eval}

The pre-award screen was run against the eight cases whose judgments state enough
to populate any profile field, filling a field only where the judgment directly
states or entails it. Two cases reach ``medium'' from bidder-side facts alone:
Braceurself, with a stated 80.25\% against 82.5\% margin, and Lancashire Care,
where the court found the authority's reasons ``too vague and generic''. Both are
among the clearest cases in the corpus where a court found the process actually
defective. That co-occurrence is worth naming and is not a validated correlation;
six cases cannot establish one.

Four of the six land ``low'' despite being real litigated challenges, and this is
a finding about data rather than a failure of the screen. Judgments record legal
reasoning; they very rarely record whether the evaluation panel was procurement-
trained or the audit trail complete, which are exactly the pre-award conditions the
rules key on. The module already states that an unpopulated profile always scores
low, and that this is a gap in the input rather than a finding. The worked
examples confirm it is the actual behaviour rather than a defensive caveat.

\section{Threats to validity}
\label{sec:threats}

\textbf{Sample size.} Every rate here is a proportion over a modest $N$. The
ablation is $n=8$ per cell with no significant result; the baseline comparison is
$n=6$. The deadlock finding is the exception, holding at 18 of 18 by architectural
necessity rather than sampling luck.

\textbf{Single model, single configuration.} Everything runs on Llama 3.1 8B at
fixed per-role temperatures. Nothing here separates ``this architecture'' from
``this model under these prompts''.

\textbf{Corpus bias.} Litigated cases are a filtered minority of all disputes, and
this corpus is further filtered toward cases well documented enough to verify. The
merits group splits 6 authority wins to 8 losses, which is far more balanced than
the original five cases were, but that balance was achieved by deliberately hunting
for authority-win cases, which is itself a selection choice.

\textbf{Detector limits.} The numeric screen is string-presence matching against
the scenario description, not a fact-checker; the premise-fabrication count on
Faraday comes from a phrase search and is a lower bound; the BATNA classifier is a
documented keyword heuristic. Each is stated in its own module.

\textbf{Structural mismatch.} Nine of 23 cases were decided at an interim stage
this system cannot model at all. Any apparent correspondence on those is
directional at best and was excluded from every accuracy figure.

\section{Reflection on the process}
\label{sec:reflection}

Three things went differently from the plan, and the differences were productive.

The first is that the interesting problem moved. The plan treated the negotiation
system as the deliverable and evaluation as a final phase. In practice the
negotiation was working within weeks and the remaining months went almost entirely
into evaluation, because the question that turned out to matter was not whether
agents could negotiate but whether anything they produced could be trusted. The
report reflects that reallocation.

The second is that the client reframed the project in one sentence. Fusion21's
interest was prevention, not replay, and the risk screen exists because of it. It
is the smallest component in the system, has no model call, and is the only one
with an observable ground truth in principle --- which makes it, awkwardly, the
most scientifically tractable thing here, and the thing this project can least
validate without data only the client can supply.

The third is that the most useful bugs came from documents I did not choose. A
board demonstration ended with a Fusion21 employee handing me an award letter and
a scoring table and asking what the system made of them. It found a genuine
arithmetic defect --- the Court agent re-applying an already-applied percentage
weighting to a sub-score, turning a correct evaluation into a false manifest-error
finding --- that three increasingly explicit prompt revisions failed to fix. The
model would state the correct approach in words and then execute the wrong
arithmetic in the same response. Those revisions were reverted rather than
shipped, since they added prompt length for no measured benefit, and the
limitation is recorded here as open. The same session surfaced a real
premise-fabrication instance which \emph{was} fixed. Choosing my own test inputs
had, without my noticing, been selecting for the failures I already knew about.
